\SourceLeaf{43}\begin{center}
{\Large SECTIO TERTIA}\par\bigskip
{\small DE}\par\bigskip
{\Large RESIDUIS POTESTATUM.}
\end{center}

\begin{center}\rule{0.09\linewidth}{0.3pt}\end{center}

\begin{center}\small\emph{Residua terminorum progressionis geometricae ab unitate incipientis constituunt seriem periodicam.}\end{center}

\Article{45}

\textsc{Theorema.} \emph{In omni progressione geometrica $1,a,aa,a^3$ etc. praeter primum $1$, alius adhuc datur terminus $a^t$, secundum modulum $p$ ad $a$ primum unitati congruus, cuius exponens $t<p$.}

\emph{Demonstr.} Quoniam modulus $p$ ad $a$, adeoque ad quamvis ipsius $a$ potestatem est primus, nullus progressionis terminus erit $\equiv0\ (\mathrm{mod.}p)$, sed quivis alicui ex his numeris $1,2,3....p-1$ congruus. Quorum multitudo quum sit $p-1$, manifestum est, si plures quam $p-1$ progressionis termini considerentur, omnes residua minima diversa habere non posse. Quocirca inter terminos $1,a,aa,a^3....a^{p-1}$ bini ad minimum congrui invenientur. Sit itaque $a^m\equiv a^n$ et $m>n$, fietque dividendo per $a^n$, $a^{m-n}\equiv1$ (art. 22) ubi $m-n<p$, et $>0$.\par
\noindent\mbox{\emph{Q. E. D.}}

\emph{Ex.} In progressione $2,4,8$ etc. terminus primus qui secundum modulum $13$ unitati est congruus, invenitur $2^{12}=4096$. At secundum modulum $23$ in eadem progressione fit $2^{11}=2048\equiv1$. Similiter numeri $5$ potestas sexta, $15625$, unitati congrua secundum modulum $7$, quinta vero, $3125$, secundum $11$. In aliis igitur casibus potestas exponentis minoris quam $p-1$ unitati congrua evadit, in aliis contra usque ad potestatem $p-1^{\mathrm{tam}}$ ascendere necesse est.

\SourceLeaf{44}\Article{46}

Quando progressio ultra terminum qui unitati est congruus continuatur, eadem quae ab initio habebantur residua prodeunt iterum. Scilicet si $a^t\equiv1$, erit $a^{t+1}\equiv a,a^{t+2}\equiv aa$ etc. donec ad terminum $a^{2t}$ perveniatur, cuius residuum minimum iterum erit $\equiv1$, atque residuorum \emph{periodum} denuo inchoat. Habetur itaque periodus $t$ residua comprehendens, quae simulac finita est ab initio semper repetitur; neque alia residua quam quae in hac periodo continentur in tota progressione occurrere possunt. Generaliter erit $a^{mt}\equiv1$, et $a^{mt+n}\equiv a^n$, id quod per designationem nostram ita exhibetur:

\[\text{\emph{Si }}r\equiv\rho\ (mod.\,t),\quad\text{\emph{erit}}\quad a^r\equiv a^\rho\ (mod.\,p).\]

\Article{47}

Petitur ex hoc theoremate compendium potestatum quantumvis magno exponente affectarum residua expedite inveniendi, simulac potestas unitati congrua innotescat. Si \emph{ex. gr.} residuum e divisione potestatis $3^{1000}$ per $13$ oriundum quaeritur, erit propter $3^3\equiv1\ (\mathrm{mod.}\ 13)$, $t\equiv3$; quare quum sit $1000\equiv1\ (\mathrm{mod.}\ 3)$, erit $3^{1000}\equiv3\ (\mathrm{mod.}\ 13)$.

\Article{48}

Quando $a^t$ est \emph{infima} potestas unitati congrua (praeter $a^0=1$, ad quem casum hic non respicimus), illi $t$ termini, residuorum periodum constituentes omnes erunt diversi, uti ex demonstratione art. 45 nullo negotio perspicitur. Tum autem propositio art. 46 converti potest; scilicet si $a^m\equiv a^n\ (\mathrm{mod.}p)$, erit $m\equiv n\ (\mathrm{mod.}t)$. Si enim $m,n$ secundum modulum $t$ incongrui essent, residua eorum minima $\mu,\nu$ diversa forent. At $a^\mu\equiv a^m,a^\nu\equiv a^n$, quare $a^\mu\equiv a^\nu$ \emph{i. e.} non omnes potestates infra $a^t$ incongruae forent contra hypoth.

Si itaque $a^k\equiv1\ (\mathrm{mod.}p)$, erit $k\equiv0\ (\mathrm{mod.}t)$ \emph{i. e.} $k$ per $t$ divisibilis.

Hactenus de modulis quibuscunque si modo ad $a$ sint primi diximus. Iam modulos qui sunt numeri absolute primi seorsim consideremus atque huic fundamento investigationem generaliorem postea superstruamus.

\SourceLeaf{45}\begin{center}\small\emph{Considerantur primo moduli qui sunt numeri primi.}\end{center}

\Article{49}

\textsc{Theorema.} \emph{Si $p$ est numerus primus ipsum $a$ non metiens, atque $a^t$ infima ipsius $a$ potestas secundum modulum $p$ unitati congrua, exponens $t$ aut erit $=p-1$ aut pars aliquota huius numeri.}

Conferantur exempla art. 45.

\emph{Demonstr.} Quum iam ostensum sit, $t$ esse aut $=p-1$, aut $<p-1$, superest, ut in posteriori casu $t$ semper ipsius $p-1$ partem aliquotam esse evincatur.

I. Colligantur residua minima positiva omnium horum terminorum $1,a,aa....a^{t-1}$, quae per $\alpha,\alpha',\alpha''$ etc. designentur, ita ut sit $\alpha=1,\alpha'\equiv a,\alpha''\equiv aa$ etc. Perspicuum est, haec omnia fore diversa, si enim duo termini $a^m,a^n$ eadem praeberent, foret (supponendo $m>n$) $a^{m-n}\equiv1$ atque $m-n<t$, \mbox{\emph{Q. E. A.}} quum nulla inferior potestas quam $a^t$ unitati sit congrua (\emph{hyp.}). Porro omnes $\alpha,\alpha',\alpha''$ etc. in serie numerorum $1,2,3...p-1$ continentur, quam tamen non exhaurient, quum $t<p-1$. Complexum omnium $\alpha,\alpha',\alpha''$ etc. per $(A)$ designabimus. Comprehendet igitur $(A)$ terminos $t$.

II. Accipiatur numerus quicunque $\beta$ ex his $1,2,3...p-1$, qui in $(A)$ desit. Multiplicetur $\beta$ per omnes $\alpha,\alpha',\alpha''$ etc., sintque residua minima inde oriunda $\beta,\beta',\beta''$ etc., quorum numerus etiam erit $t$. At haec residua tum inter se quam ab omnibus $\alpha,\alpha',\alpha''$ etc. erunt diversa. Si enim \emph{prior} assertio falsa esset, haberetur $\beta a^m\equiv\beta a^n$ adeoque dividendo per $\beta$, $a^m\equiv a^n$, contra ea quae modo demonstravimus: si vero \emph{posterior}, haberetur $\beta a^m\equiv a^n$, unde, quando $m<n$, $\beta\equiv a^{n-m}$ \emph{i. e.} $\beta$ alicui ex his $\alpha,\alpha',\alpha''$ etc. congruus contra hyp.; quando vero $m>n$, sequitur multiplicando per $a^{t-m}$, $\beta a^t\equiv a^{t+n-m}$, sive propter $a^t\equiv1$, $\beta\equiv a^{t+n-m}$, quae est eadem absurditas. Designetur complexus omnium $\beta,\beta',\beta''$ etc. quorum multitudo $=t$, per $(B)$, habebunturque iam $2t$ numeri ex his $1,2,3...p-1$. Quodsi igitur $(A)$ et $(B)$ omnes hos numeros complectuntur, fit $\frac{p-1}{2}=t$ adeoque theorema demonstratum.

III. Si vero aliqui adhuc deficiunt, sit horum aliquis $\gamma$. Per hunc multiplicentur omnes $\alpha,\alpha',\alpha''$ etc., productorumque residua minima sint $\gamma,\gamma',\gamma''$ etc.; omnium complexus per $(C$ designetur. $(C)$ igitur comprehendet $t$ numeros ex his $1,2,3...p-1$, qui omnes tum inter se quam a numeris in $(A)$ et $(B)$
\SourceLeaf{46}contentis erunt diversi. Assertiones priores eodem modo demonstrantur ut in II, tertia ita. Si esset $\gamma a^m\equiv\beta a^n$, fieret $\gamma\equiv\beta a^{n-m}$, aut $\equiv\beta a^{t+n-m}$ prout $m<n$, aut $>n$, in utroque casu $\gamma$ alicui ex $(B)$ congrua contra hyp. Habentur igitur $3t$ numeri ex his $1,2,3...p-1$, atque si nulli amplius desunt, fiet $t=\frac{p-1}{3}$ adeoque theorema erit demonstratum.

IV. Si vero etiamnum aliqui desunt, eodem modo ad quartum numerorum complexum $(D)$ progrediendum erit etc. Patet vero quoniam numerorum $1,2,3....p-1$ multitudo est finita, tandem eam exhaustum iri, adeoque multiplum ipsius $t$ fore: quare $t$ erit pars aliquota numeri $p-1$. \mbox{\emph{Q. E. D.}}

\begin{center}\small\emph{Fermatii Theorema.}\end{center}

\Article{50}

Quum igitur $\frac{p-1}{t}$ sit integer, sequitur evehendo utramque partem congruentiae $a^t\equiv1$ ad potestatem exponentis $\frac{p-1}{t}$, $a^{p-1}\equiv1$, \emph{sive $a^{p-1}-1$ semper per $p$ divisibilis est, quando $p$ est primus ipsum $a$ non metiens.}

Theorema hoc quod tum propter elegantiam tum propter eximiam utilitatem omni attentione dignum, ab inventore \emph{theorema Fermatianum} appellari solet. Vid. \emph{Fermatii Opera Mathem. Tolosae} 1679 \emph{fol. p.} 163. Demonstrationem inventor non adiecit, quam tamen in potestate sua esse professus est. Ill. Euler primus demonstrationem publici iuris fecit, in diss. cui titulus \emph{Theorematum quorundam ad numeros primos spectantium demonstratio, Comm. Acad. Petrop. T.} VIII\begingroup\renewcommand{\thefootnote}{*)}\footnote{In comment. anteriore vir summus ad scopum nondum pervenerat. \emph{Comm. Petr. T.} VI \emph{p.} 106. --- In controversia famosa inter Maupertuis et Konig, a principio actionis minimae orta, sed mox ad res heterogeneas egressa, K\"onig in manibus se habere dixit autographum Leibnitianum, in quo demonstratio huius theorematis cum Euleriana prorsus conspirans contineatur. \emph{Appel au public. p.} $106$. Licet vero fidem huic testimonio denegare nolimus, certe Leibnitius inventum suum numquam publicavit. Conf. \emph{Hist. de l'Ac. de Prusse, A.} 1750 \emph{p.} 530.}\endgroup. Innititur ista evolutioni potestatis $(a+1)^p$, ubi ex coefficientium forma facillime deducitur $(a+1)^p-a^p-1$ semper per $p$ fore divisibilem, adeoque $(a+1)^p-(a+1)$ per $p$ divisibilem fore, quando $a^p-a$ per $p$ sit divisibilis. Iam quia $1^p-1$ semper per $p$ divisibilis est, etiam $2^p-2$ semper erit; hinc etiam $3^p-3$ etc. generaliterque $a^p-a$. Quodsi itaque $p$ ipsum $a$ non metitur, etiam $a^{p-1}-1$ per $p$ divisibilis erit. Haec sufficient ad methodi indolem declarandam. Clar. Lambert similem demonstrationem tradidit in \emph{Actis Erudit.} 1769
